前三部分讨论了构建高级大语言模型智能体的 AI 技术——这些智能体能够与现实世界交互并在其中行动。这些能力使它们能够为用户做有用的事。在本部分中，我们将讨论大语言模型如何使用工具来执行专门任务，例如查询信息或控制机器人。智能体通过学到的 API（应用程序编程接口）来使用工具。

工具可以帮助我们解决预训练大语言模型的一个主要问题：模型中所包含的信息，其时效性仅止于训练所用的数据集。检索增强技术使大语言模型能够提供最新的信息。  

这类智能体已经出现了许多种，在本部分最后一章中，我们将回顾来自医疗、金融与科学领域的重要实例。最后，我们将讨论一个实践案例研究，把各项技术汇聚起来——我们将实现一个大语言模型智能体来完成医疗诊断。
